\section{Analyse des résultats, complement}
\label{sec:analyse}

\subsection{DMA}

Bien que nous aillons traite les donnees en Sec. \ref{sec:DMAresults}, le traitement permet de determiner plus precisement le pic de chacune des courbes en Fig. \ref{fig:DMAresults}. Couple a une autre fonction logiciel (recherche des extremeums locaux lors de changement de variation d'une fonction continue). Nous obtenons en sortie du programme la figure suivante :
\\
\begin{table}[!ht]
\centering
\footnotesize
\setlength{\tabcolsep}{4pt}
\renewcommand{\arraystretch}{0.9}
\begin{tabular}{|c|c|c|}
\hline
	\textbf{Frequence [Hz]} 
	& \textbf{Peak [$^\circ$]}
	& \textbf{incertitude [$^\circ$]}
\\\hline
	1.000 & 140.5 & 1.606 \\
	2.000 & 143.5 & 1.594 \\
	5.000 & 148.6 & 1.601 \\
	10.000 & 150.5 & 1.616 \\
	20.000 & 155.7 & 1.518 \\
\hline
\end{tabular}
	\caption{Valeurs de temperatures auxquelles E (module d'Young) est maximal}
\label{tab:peaks}
\end{table}


On peut donc extraire tous les extremums locaux de Fig. \ref{fig:DMApeaks} que l'on peut retrouver en annexe \ref{annexe:peaks}. Si l'on plot tout ca nous obtenons le graph d'Arrenhius, soit : 

\begin{figure}[hct]
        \centering
        \includegraphics[scale=0.5]{Seite_arrhenius_graph.pdf}
        \label{fig:arrenius}
        \caption{Arrhenius graph}
\end{figure}

\newpage

\subsection{DSC}

Dans son ensemble chaque experience en DSC est visible en Fig. \ref{fig:combined}.

\begin{figure}[hct]
        \centering
        \includegraphics[scale=0.5]{combined.pdf}
        \label{fig:arrenius}
        \caption{Arrhenius graph}
\end{figure}

\\
On peut concentrer les donnees interessantes dans le tableau suivant :
\\
\begin{table}[!ht]
\centering
\footnotesize
\setlength{\tabcolsep}{4pt}
\renewcommand{\arraystretch}{0.9}
\begin{tabular}{|c|c|c|c|c|}
\hline
	\textbf{Variation}
	& \textbf{$V_v_t$ [$K/min$]}
        & \textbf{Pic complexe [$^\circ$]}
        & \textbf{Aire [$J/g$]}
	& \textbf{$T_v$ [$^\circ$]}
\\\hline
	Chauffe & 50.0 & 244.8 & -43.15 & Oui  \\
	Refroidissement & -200.0 & $\emptyset$ & $\emptyset$ & $\emptyset$  \\
	Chauffe & 10.0 & 131.1-245.9 & 23.93-(-35.58) & Oui \\
	Refroidissment & -10.0 & 186.8 & 36.19 & $\emptyset$ \\
\hline
\end{tabular}
        \caption{$V_v_t$: vitesse de la variation thermique, $T_v$: Temperature de transition vitreuse}
\label{val:peaks}
\end{table}
